\section{Radio-Performance Distributions}
\label{sec:results_sinr}

This section reports radio-performance statistics for the evaluated placement strategies. The panels summarize the empirical \gls{sinr} distribution across user snapshots and the aggregate \gls{esr} distribution across evaluated time samples.

\begin{figure*}[t]
    \centering

    \ref{sharedlegend}
    \vspace*{1em}
    
    \begin{subfigure}[t]{0.48\textwidth}
        \centering
        \begingroup
        \def\postprocessfigurewidth{\linewidth}
        \pgfplotsset{table/search path={\SINRFigureRoot}}
        \input{\SINRCdfFigurePath}
        \endgroup
        \caption{Empirical \gls{sinr} distribution across \num{7210} user snapshots: \num{10} scheduled users over \num{721} time samples.}
        \label{fig:sinr_cdf}
    \end{subfigure}\hfill
    \begin{subfigure}[t]{0.48\textwidth}
        \centering
        \begingroup
        \def\postprocessfigurewidth{\linewidth}
        \pgfplotsset{table/search path={\SINRFigureRoot}}
        \input{\ESRCdfFigurePath}
        \endgroup
        \caption{\gls{esr} \gls{cdf} across the \num{721} evaluated time samples.}
        \label{fig:esr_cdf}
    \end{subfigure}
    \caption{Radio-performance distributions for the campus trace. \strategy4[title] produces the most favorable \gls{sinr} and \gls{esr} distributions, with \strategy5[title] as the next-best movable strategy.}
    \label{fig:radio_primary_pair}
\end{figure*}

\Cref{fig:radio_primary_pair} establishes a clear ordering. The nearest-user proxy-\gls{esr} controller of \strategy4[title] dominates almost the full empirical \gls{sinr} range. Its mean \gls{sinr} reaches \SI{39.7}{dB}, with a \num{10}th percentile of \SI{24.5}{dB}; \strategy2[title] reaches \SI{21.3}{dB} mean and \SI{7.5}{dB} at the \num{10}th percentile. The outage fraction below \SI{0}{dB} drops from \SI{4.41}{\percent} for \strategy2[title] to \SI{1.87}{\percent} for \strategy4[title]. \strategy5[title] is the second-best radio strategy with \SI{35.7}{dB} mean \gls{sinr} and \SI{16.3}{dB} \num{10}th-percentile \gls{sinr}, but its outage fraction remains essentially identical to \strategy2[title] at about \SI{4.42}{\percent}. Hence, the power-only proxy improves the middle and upper portions of the distribution much more strongly than the weakest-link tail.

\strategy1[title] is weakest throughout, with \SI{9.8}{dB} mean \gls{sinr}, \SI{-22.1}{dB} \num{10}th-percentile \gls{sinr}, and \SI{18.4}{\percent} outage.

The \gls{esr} panel in \cref{fig:radio_primary_pair} follows the same ranking. Mean \gls{esr} reaches \SI{133.1}{bit/s/Hz} for \strategy4[title], \SI{122.8}{bit/s/Hz} for \strategy5[title], \SI{74.9}{bit/s/Hz} for \strategy2[title], and \SI{56.4}{bit/s/Hz} for \strategy1[title]. Relative to \strategy2[title], \strategy4[title] improves the mean \gls{esr} by \SI{77.8}{\percent} and the \num{10}th-percentile \gls{esr} by \SI{88.8}{\percent}. \strategy5[title] still improves the mean \gls{esr} by \SI{64.1}{\percent}, but remains consistently below \strategy4[title].

Together, these figures indicate that the best result in this scenario is not merely ``more movement'' or ``more power'', but the combination of relocation with a phase-aware \gls{ue}-to-\gls{ue} proxy-\gls{csi} objective. \strategy2[title] remains a strong zero-relocation benchmark, whereas \strategy1[title] is clearly mismatched to the spatially spread campus user process.
% \FloatBarrier
